\documentclass[a4paper,12pt]{report}
\usepackage{ctex}
\usepackage{xcolor}
\usepackage{graphicx}
\usepackage{caption}
\usepackage{listings}
\usepackage{geometry}
\usepackage{hyperref}

\geometry{left=2.8cm,right=2.8cm,top=2.7cm,bottom=2.7cm}
\setlength{\parindent}{2em}
\setlength{\parskip}{0.3em}
\setlength{\emergencystretch}{2em} 

\definecolor{codegray}{RGB}{247,247,247}
\definecolor{framegray}{RGB}{150,150,150}
\lstset{
  basicstyle=\ttfamily\small,
  backgroundcolor=\color{codegray},
  frame=single,
  rulecolor=\color{framegray},
  breaklines=true,
  keepspaces=true,
  columns=fullflexible, 
  showstringspaces=false
} 

\newcommand{\reportimage}[2]{%
  \begin{center}
    \includegraphics[width=0.88\textwidth,height=0.68\textheight,keepaspectratio]{#1}
    \captionof{figure}{#2}
  \end{center}  
}

\begin{document}
\title{系统开发工具基础}
\author{肖扬\\计算机科学与技术\\学号：25020007140}
\date{2026年9月14日}
\maketitle

\pagenumbering{roman}
\tableofcontents
\newpage
\pagenumbering{arabic}

\color{red}\chapter{实验环境说明}\color{black}
\begin{enumerate}
  \item {\large 系统环境：Ubuntu，用户名 xy\_ai，WSL2 架构}
  \item {\large Shell：bash 5.1.16（x86\_64-pc-linux-gnu）}
  \item {\large Python：Miniconda 环境中的 Python 3.12；第 12 题使用安装了 PyTorch 的 AI\_env 环境进行 CPU 训练}
  \item {\large 构建工具：build 1.6.1、setuptools 68 及以上版本}
  \item {\large LaTeX：TeX Live 2022/dev（Debian 版），使用 XeLaTeX 编译}
\end{enumerate}

\color{red}\chapter{课堂学习部分}\color{black}

\section{第 9 题 从源码构建并在干净环境安装 Wheel}

\reportimage{images/q09-01-task.png}{第 9 题题目要求}

\begin{enumerate}
  \item {\large 创建项目目录和所需文件}

  在 assignment-3 下创建 \texttt{q09}，再建立 \texttt{src/greetlab} 目录，并创建
  \texttt{\_\_init\_\_.py}、\texttt{cli.py} 和 \texttt{pyproject.toml}。最后用
  \texttt{find} 查看目录结构，文件均已放到指定位置。

\reportimage{images/q09-02-create-files.png}{创建 q09 项目目录及文件}

  \item {\large 编写命令行程序}

  \texttt{src/greetlab/cli.py} 的内容如下：

\begin{lstlisting}[language=Python]
import argparse

def main():
    p = argparse.ArgumentParser()
    p.add_argument("--name", required=True)
    a = p.parse_args()
    print(f"Hello, {a.name}!")
\end{lstlisting}

  程序使用 \texttt{argparse} 接收必填的 \texttt{--name} 参数，然后输出问候语。

  \item {\large 编写项目配置文件}

  \texttt{pyproject.toml} 的内容如下，项目名末尾加上了学号：

\begin{lstlisting}
[build-system]
requires = ["setuptools>=68"]
build-backend = "setuptools.build_meta"

[project]
name = "greetlab-25020007140"
version = "0.1.0"

[project.scripts]
sdt-greet = "greetlab.cli:main"
\end{lstlisting}

  \item {\large 安装 build 工具并构建 Wheel}

  第一次执行 \texttt{python -m build} 时，终端提示缺少 \texttt{build} 模块。我先运行
  \texttt{python -m pip install build} 完成安装，再重新构建。随后程序创建隔离环境，读取
  \texttt{pyproject.toml} 并生成发布文件。

\reportimage{images/q09-03-build.png}{安装 build 工具并执行项目构建}

  构建完成后，\texttt{dist} 目录中生成以下两个文件：

\begin{lstlisting}
greetlab_25020007140-0.1.0-py3-none-any.whl
greetlab_25020007140-0.1.0.tar.gz
\end{lstlisting}

  \item {\large 在干净虚拟环境中安装 Wheel}

  返回 assignment-3 目录，使用下面的命令新建并启用虚拟环境：

\begin{lstlisting}[language=bash]
python -m venv q09-test
source q09-test/bin/activate
pip install q09/dist/greetlab_25020007140-0.1.0-py3-none-any.whl
\end{lstlisting}

  安装命令只指定刚生成的 Wheel。终端显示
  \texttt{Successfully installed greetlab-25020007140-0.1.0}，安装顺利完成。

  \item {\large 在 q09 目录外验证命令行程序}

  仍在 assignment-3 目录中运行：

\begin{lstlisting}[language=bash]
sdt-greet --name 25020007140
\end{lstlisting}

  终端输出 \texttt{Hello, 25020007140!}。入口命令在 \texttt{q09} 目录外也能运行，
  因此不是依靠当前源码目录才生效。完成后使用 \texttt{deactivate} 退出虚拟环境。

\reportimage{images/q09-04-install-test.png}{在干净虚拟环境中安装并验证 Wheel}
\end{enumerate}

\section{第 10 题 让编程智能体进入可验证的修复循环}

\reportimage{images/q10-01-task.png}{第 10 题题目要求}

\begin{enumerate}
  \item {\large 复制第 9 题项目并编写测试}

  在 assignment-3 目录下把 \texttt{q09} 复制为 \texttt{q10}，再建立
  \texttt{tests/test\_cli.py}，用于检查空白用户名的处理结果。

\begin{lstlisting}[language=bash]
cp -r q09 q10
cd q10
mkdir -p tests
nano tests/test_cli.py
\end{lstlisting}

  测试文件如下。这里用 \texttt{monkeypatch} 把 \texttt{--name} 设置为三个空格，
  再检查 \texttt{main()} 是否抛出 \texttt{SystemExit}，并核对退出状态码是否为 2。

\begin{lstlisting}[language=Python]
import sys
import pytest

from greetlab.cli import main


def test_blank_name_exits(monkeypatch):
    monkeypatch.setattr(sys, "argv", ["sdt-greet", "--name", "   "])
    with pytest.raises(SystemExit) as exc:
        main()

    assert exc.value.code == 2
\end{lstlisting}

  \item {\large 运行测试并确认失败}

  在 \texttt{base} 环境的 \texttt{q10} 目录中运行：

\begin{lstlisting}[language=bash]
python -m pytest -q
\end{lstlisting}

  终端显示 \texttt{1 failed}，报错为 \texttt{DID NOT RAISE SystemExit}。
  程序仍输出了空白问候语，说明它把三个空格当成了正常用户名。复现问题后，
  我再把修复任务交给编程智能体。

\reportimage{images/q10-redo-01-failed.png}{创建测试并确认空白用户名测试失败}

  \item {\large 向编程智能体说明目标和约束}

  在 Codex 中说明当前项目的位置，并给出修复目标、允许修改的文件和验证命令。
  核心提示如下：

  \begin{quote}
  当 \texttt{--name} 只包含空白字符时，让 \texttt{main} 以
  \texttt{SystemExit(2)} 结束。只修改 \texttt{q10/src/greetlab/cli.py}，
  不要修改测试，保留正常用户名的输出。请先检查失败原因，修复后运行
  \texttt{python -m pytest -q}。
  \end{quote}

  智能体读取代码和测试后先运行测试，得到相同的失败结果。原因是 \texttt{argparse}
  会接收三个空格作为参数值，而原程序解析参数后直接打印，没有继续检查字符串内容。

\reportimage{images/q10-redo-02-agent-process.png}{向智能体说明任务并由其检查、修复和验证}

  \item {\large 检查智能体的修复结果}

  智能体在参数解析后增加以下判断，测试文件保持不变：

\begin{lstlisting}[language=Python]
a = p.parse_args()

if not a.name.strip():
    raise SystemExit(2)

print(f"Hello, {a.name}!")
\end{lstlisting}

  \texttt{strip()} 用来判断用户名去掉两端空格后是否为空。若为空就抛出
  \texttt{SystemExit(2)}；否则仍使用原来的 \texttt{a.name} 输出问候语。
  智能体修改后执行 \texttt{python -m pytest -q}，结果为 \texttt{1 passed}，
  并检查了正常用户名的输出。

\reportimage{images/q10-redo-03-agent-result.png}{智能体说明修改内容及测试结果}

  \item {\large 人工检查差异并再次运行测试}

  回到 \texttt{q10} 目录，手动执行下面的命令，对照第 9 题的源码检查修改，
  然后再次运行测试：

\begin{lstlisting}[language=bash]
diff -u ../q09/src/greetlab/cli.py src/greetlab/cli.py
python -m pytest -q
\end{lstlisting}

  差异中只有空白用户名判断、退出语句和用于分隔代码的空行，没有无关修改。
  我再次运行测试，终端显示 \texttt{1 passed in 0.01s}。

\reportimage{images/q10-redo-04-human-review.png}{人工检查源码差异并再次运行测试}

  \newpage
  \item {\large 记录核心提示和验证结果}

  在 \texttt{q10/ai\_log.md} 中用四行记录本次任务的核心提示、测试命令、
  智能体改动和人工验证结果，没有超过题目要求的 5 行。

\begin{lstlisting}
核心提示：空白 name 以 SystemExit(2) 结束；只改 cli.py，不改测试，保留正常输出。
测试命令：在 q10 目录运行 python -m pytest -q。
智能体改动：解析参数后增加 strip() 空白检查；测试由 1 failed 变为 1 passed。
人工验证：检查 diff，无无关修改；再次运行测试，结果为 1 passed。
\end{lstlisting}
\end{enumerate}

\section{第 11 题 把“无法处理”的协作材料改成可执行信息}

\reportimage{images/q11-01-task.png}{第 11 题题目要求}

\begin{enumerate}
  \item {\large 创建 communication.md}

  在 assignment-3 下建立 \texttt{q11} 目录，进入目录后创建
  \texttt{communication.md}，用于填写改写后的协作材料。

\begin{lstlisting}[language=bash]
mkdir q11
cd q11
nano communication.md
\end{lstlisting}

\reportimage{images/q11-02-create-file.png}{创建 q11 目录和 communication.md}

  \item {\large 改写 Issue}

  Issue 中写明运行环境为 Windows，给出复现命令
  \texttt{sdt-greet --name "   "}。期望结果是拒绝只含空白字符的 name，并以非零状态退出；
  实际结果是程序输出空白问候语并以状态码 0 退出。不能确定的问题写为“是否仅在 Windows 下出现，待确认”。

  \item {\large 改写提交信息}

  提交标题使用祈使语气：

\begin{lstlisting}
fix: reject blank name input
\end{lstlisting}

  正文说明处理方法：拒绝只含空白字符的用户名，并使用 \texttt{SystemExit(2)} 退出。

  \item {\large 改写评审意见}

  评审中用 \texttt{Blocking} 指出程序把无效输入当成了有效输入，并建议通过
  \texttt{strip()} 检查空白；再用 \texttt{Suggestion} 建议补充自动化测试，验证空白输入会触发
  \texttt{SystemExit(2)}。两条意见分别对应必须修复的问题和后续改进建议。

  \item {\large 检查内容和长度}

  使用 \texttt{wc -m communication.md} 统计字符数，结果为 374，没有超过题目规定的
  400 字。随后用 \texttt{cat communication.md} 查看全文，三部分内容均已写入。

\reportimage{images/q11-03-check-content.png}{检查 communication.md 的字符数和内容}
\end{enumerate}

\section{第 12 题 修复一个可复现的线性回归训练循环}

\reportimage{images/q12-01-task.png}{第 12 题题目要求}

\begin{enumerate}
  \item {\large 激活环境并创建训练文件}

  本题需要使用 PyTorch，因此激活已经安装 PyTorch 的 \texttt{AI\_env} 环境，进入
  assignment-3 下的 \texttt{q12} 目录并创建 \texttt{train.py}。

\begin{lstlisting}[language=bash]
conda activate AI_env
cd ~/系统工具开发基础/assignment-3/q12
nano train.py
\end{lstlisting}

  \item {\large 补全训练循环}

  根据题目给出的代码，在每轮训练中依次计算预测值和损失，再调用
  \texttt{zero\_grad()}、\texttt{backward()} 和 \texttt{step()}：

\begin{lstlisting}[language=Python]
for _ in range(200):
    pred = model(x)
    loss = loss_fn(pred, y)

    opt.zero_grad()
    loss.backward()
    opt.step()
\end{lstlisting}

  \item {\large 计算并输出最终结果}

  训练结束后调用 \texttt{model.eval()}，并在 \texttt{torch.no\_grad()} 中计算最终损失。
  最后打印 final loss、weight 和 bias。代码没有把张量和模型移到 CUDA，所以本次训练使用 CPU。

\reportimage{images/q12-02-code.png}{补全后的 train.py 代码}

  \item {\large 运行程序并检查结果}

  连续两次执行 \texttt{python train.py}，得到相同结果：
  \begin{itemize}
    \item final loss：\texttt{2.3923129102709773e-12}，小于 0.001；
    \item weight：\texttt{2.99999737739563}，接近 3；
    \item bias：\texttt{-0.9999998807907104}，接近 -1。
  \end{itemize}

\reportimage{images/q12-03-result.png}{两次运行线性回归训练程序的结果}
\end{enumerate}

\color{red}\chapter{课后习题}\color{black}

\section{第 1 题 创建并安装带有 pyproject.toml 的 Python 包}

\begin{enumerate}
  \item {\large 创建项目配置与包目录}

  在 \texttt{assignment-3/homework/01} 目录中创建 \texttt{pyproject.toml}，使用
  setuptools 作为构建后端，并填写软件包名称、版本、说明和 Python 版本要求：

\begin{lstlisting}
[build-system]
requires = ["setuptools>=68", "wheel"]
build-backend = "setuptools.build_meta"

[project]
name = "hello-package"
version = "0.1.0"
description = "A simple Python package"
requires-python = ">=3.10"
\end{lstlisting}

  接着创建 \texttt{hello\_package} 目录和 \texttt{\_\_init\_\_.py}，得到一个最小的
  Python 包结构。

\begin{lstlisting}[language=bash]
mkdir hello_package
touch hello_package/__init__.py
\end{lstlisting}

  \item {\large 创建虚拟环境并安装软件包}

  使用 Python 自带的 venv 模块创建并激活虚拟环境。\texttt{which python} 返回当前项目下的
  \texttt{.venv/bin/python}，后面的安装操作都在这个环境中完成。

\begin{lstlisting}[language=bash]
python3 -m venv .venv
source .venv/bin/activate
which python
pip install .
pip list
\end{lstlisting}

  pip 根据 \texttt{pyproject.toml} 构建 Wheel 并完成安装。随后查看软件包列表，
  其中 \texttt{hello-package} 的版本为 0.1.0。

\noindent{\color{red}操作失误，第三次实验作业不小心放到第四次里面了，现在已经移动到第三次实验了。}

\reportimage{images/hw01-01-package-install.png}{创建 Python 包并在虚拟环境中安装}

  \item {\large 创建并检查锁文件}

  在虚拟环境中安装 uv，查看版本后运行 \texttt{uv lock}，生成 \texttt{uv.lock}：

\begin{lstlisting}[language=bash]
pip install uv
uv --version
uv lock
cat uv.lock
head -n 40 uv.lock
\end{lstlisting}

  本次使用的 uv 版本为 0.12.15。锁文件开头写有格式版本、修订号和
  \texttt{requires-python = \">=3.10\"}；软件包条目记录了 \texttt{hello-package}、
  版本 0.1.0 和本地可编辑来源，与 \texttt{pyproject.toml} 中的配置一致。

\noindent{\color{red}操作失误，第三次实验作业不小心放到第四次里面了，现在已经移动到第三次实验了。}

\reportimage{images/hw01-02-lock-file.png}{使用 uv 生成并检查锁文件}
\end{enumerate}

\newpage
\section{第 2 题 使用现有的 pyproject.toml 项目执行构建并检查发布文件}

\begin{enumerate}
  \item {\large 创建并配置待构建项目}

  在 \texttt{assignment-3/homework} 下创建 \texttt{02} 目录，再建立
  \texttt{hello\_build} 包目录和 \texttt{\_\_init\_\_.py}。现有的
  \texttt{pyproject.toml} 内容如下：

\begin{lstlisting}
[build-system]
requires = ["setuptools"]
build-backend = "setuptools.build_meta"

[project]
name = "hello-build"
version = "0.1.0"
description = "A simple build example"
requires-python = ">=3.10"
\end{lstlisting}

\begin{lstlisting}[language=bash]
mkdir 02
cd 02
mkdir hello_build
touch hello_build/__init__.py
\end{lstlisting}

  \item {\large 执行构建并处理依赖获取失败}

  第一次构建时，隔离环境需要从 PyPI 获取 \texttt{setuptools} 和 \texttt{wheel}，
  但当时网络连接被拒绝，终端提示 \texttt{Failed to build}。调整
  \texttt{pyproject.toml} 后，我加上 \texttt{--no-isolation}，直接使用当前 Python 环境中
  已安装的构建依赖：

\begin{lstlisting}[language=bash]
python3 -m build --no-isolation
\end{lstlisting}

  这次构建顺利完成，同时生成了源码分发包和 Wheel。

\noindent{\color{red}操作失误，第三次实验作业不小心放到第四次里面了，现在已经移动到第三次实验了。}

\reportimage{images/hw02-01-build-process.png}{创建项目并处理隔离构建的依赖获取失败}

  \item {\large 确认构建工具并重复构建}

  \texttt{python3 -m pip show build} 显示 build 的版本为 1.6.1。随后再次执行无隔离构建，
  两个发布文件仍能正常生成。

\noindent{\color{red}操作失误，第三次实验作业不小心放到第四次里面了，现在已经移动到第三次实验了。}

\reportimage{images/hw02-02-build-success.png}{检查 build 版本并再次成功构建}

  \item {\large 检查生成的发布文件}

  使用 \texttt{ls dist} 查看发布目录，其中包含两个文件：

\begin{lstlisting}
hello_build-0.1.0-py3-none-any.whl
hello_build-0.1.0.tar.gz
\end{lstlisting}

  再运行 \texttt{file dist/*.whl}，终端把 Wheel 识别为使用 deflate 压缩的
  Zip archive data。也就是说，Wheel 文件可以按 ZIP 归档查看。

\noindent{\color{red}操作失误，第三次实验作业不小心放到第四次里面了，现在已经移动到第三次实验了。}

\reportimage{images/hw02-03-dist-check.png}{检查 dist 目录及 Wheel 文件类型}
\end{enumerate}

\newpage
\section{第 3 题 安装第二题构建的 Wheel 包并检查安装结果}

\begin{enumerate}
  \item {\large 创建并激活虚拟环境}

  在 \texttt{assignment-3/homework} 下创建 \texttt{03} 目录，并在其中建立独立的
  Python 虚拟环境，避免这次安装影响其他环境。

\begin{lstlisting}[language=bash]
mkdir -p 03
cd 03
python3 -m venv .venv
source .venv/bin/activate
\end{lstlisting}

  \item {\large 安装第二题生成的 Wheel}

  使用相对路径指定第二题 \texttt{dist} 目录中的 Wheel 文件并执行安装：

\begin{lstlisting}[language=bash]
pip install ../02/dist/hello_build-0.1.0-py3-none-any.whl
\end{lstlisting}

  pip 读取本地 Wheel 后安装 \texttt{hello-build}，终端显示
  \texttt{Successfully installed hello-build-0.1.0}。

  \item {\large 使用 pip show 检查安装结果}

  \texttt{pip list} 中已经出现版本为 0.1.0 的 \texttt{hello-build}。继续运行
  \texttt{pip show hello-build}，名称、版本和说明都与第二题的 \texttt{pyproject.toml}
  一致，安装位置在当前虚拟环境的 \texttt{site-packages} 目录。

\noindent{\color{red}操作失误，第三次实验作业不小心放到第四次里面了，现在已经移动到第三次实验了。}

\reportimage{images/hw03-install-wheel.png}{安装第二题构建的 Wheel 并检查安装结果}
\end{enumerate}

\newpage
\section{第 4 题 查看 Python 包的依赖信息和安装位置}

\begin{enumerate}
  \item {\large 查看软件包元数据与依赖信息}

  激活第 3 题创建的虚拟环境，再使用 \texttt{pip show} 查看已安装软件包的详细信息：

\begin{lstlisting}[language=bash]
source .venv/bin/activate
pip show hello-build
\end{lstlisting}

  输出中 \texttt{Name} 为 \texttt{hello-build}，\texttt{Version} 为 0.1.0，
  \texttt{Summary} 为 \texttt{A simple build example}。\texttt{Requires} 为空，表示该包
  没有声明运行依赖；\texttt{Required-by} 为空，表示当前环境中没有其他包依赖它。

  \item {\large 查看安装位置}

  \texttt{pip show} 的 \texttt{Location} 字段显示软件包安装在当前虚拟环境的
  \texttt{.venv/lib/python3.12/site-packages} 目录。随后导入 \texttt{hello\_build} 并输出
  模块的 \texttt{\_\_file\_\_} 属性，以确认实际加载的文件位置：

\begin{lstlisting}[language=bash,basicstyle=\ttfamily\footnotesize]
python -c "import hello_build; print(hello_build.__file__)"
\end{lstlisting}

  终端返回当前虚拟环境中 \texttt{hello\_build/\_\_init\_\_.py} 的完整路径，
  与 \texttt{pip show} 给出的安装目录相符。

\noindent{\color{red}操作失误，第三次实验作业不小心放到第四次里面了，现在已经移动到第三次实验了。}

\reportimage{images/hw04-package-info.png}{查看 hello-build 的依赖信息和实际安装位置}
\end{enumerate}

\newpage
\section{第 5 题 使用 Claude 智能体读取代码并自动生成说明文件}

\begin{enumerate}
  \item {\large 创建并运行示例程序}

  在 \texttt{assignment-3/homework} 下创建 \texttt{05} 目录，并编写
  \texttt{calculator.py}。其中 \texttt{add(a, b)} 负责加法，\texttt{multiply(a, b)}
  负责乘法。运行程序后，终端依次输出 5 和 20。

\begin{lstlisting}[language=bash]
mkdir -p 05
cd 05
nano calculator.py
python3 calculator.py
\end{lstlisting}

  \item {\large 启动 Claude 并说明任务要求}

  在项目目录运行 \texttt{claude}，进入 Claude Code 交互界面，本次版本为 2.1.236。
  我要求智能体读取 \texttt{calculator.py}，说明函数和运行结果，不改源代码，
  只在当前目录新建 \texttt{README.md}。

\noindent{\color{red}操作失误，第三次实验作业不小心放到第四次里面了，现在已经移动到第三次实验了。}

\reportimage{images/hw05-01-run-claude.png}{运行示例程序并向 Claude 智能体提交任务}

  \item {\large 由智能体分析代码并生成 README}

  Claude 读取 \texttt{calculator.py} 后找到了 \texttt{add} 和 \texttt{multiply} 两个函数，
  并写出一份 21 行的 \texttt{README.md}。文件分为程序功能、函数说明和运行结果三部分，
  其中也写明了示例输出 5 和 20。

  写完后，智能体重新读取 \texttt{README.md} 并汇总结果。\texttt{calculator.py}
  没有发生改动，符合本题要求。

\noindent{\color{red}操作失误，第三次实验作业不小心放到第四次里面了，现在已经移动到第三次实验了。}

\reportimage{images/hw05-02-agent-result.png}{Claude 分析代码并生成 README.md}

  \item {\large 退出智能体并检查生成文件}

  输入 \texttt{/exit} 退出 Claude Code，再运行 \texttt{ls}。目录中保留了
  \texttt{calculator.py}，同时多出了新生成的 \texttt{README.md}。

\noindent{\color{red}操作失误，第三次实验作业不小心放到第四次里面了，现在已经移动到第三次实验了。}

\reportimage{images/hw05-03-exit-list.png}{退出 Claude Code 并查看生成文件}

  最后用 \texttt{cat README.md} 查看文件。内容包括加法、乘法两个函数的作用、参数和返回值，
  也列出了程序运行时的输出 5 和 20。

\noindent{\color{red}操作失误，第三次实验作业不小心放到第四次里面了，现在已经移动到第三次实验了。}

\reportimage{images/hw05-04-readme.png}{检查 Claude 自动生成的 README.md}
\end{enumerate}

\newpage
\section{第 6 题 让 Claude 智能体发现并修复 Python 程序中的错误}

\begin{enumerate}
  \item {\large 创建错误程序并复现问题}

  在 \texttt{assignment-3/homework} 下创建 \texttt{06} 目录，并编写
  \texttt{divide.py}。\texttt{divide(a, b)} 直接返回 \texttt{a / b}，程序分别用 2 和 0
  作除数调用该函数。

\begin{lstlisting}[language=Python]
def divide(a, b):
    return a / b


print(divide(10, 2))
print(divide(10, 0))
\end{lstlisting}

  第一个调用输出 5.0，第二个调用则触发
  \texttt{ZeroDivisionError: division by zero}，程序在这里中断。

\noindent{\color{red}操作失误，第三次实验作业不小心放到第四次里面了，现在已经移动到第三次实验了。}

\reportimage{images/hw06-01-error.png}{运行 divide.py 并复现除零错误}

  \item {\large 向 Claude 提交自动修复任务}

  在当前目录启动 Claude Code。我要求智能体读取并运行 \texttt{divide.py}，找到报错原因后
  直接修改文件；修改完成还要再次运行，直到程序能够正常结束。

  Claude 很快定位到 \texttt{divide(10, 0)}。不过第一次运行时，Windows 环境中找不到
  \texttt{python3}；改用 WSL 后又遇到路径转换问题。调整命令后，程序终于在 Ubuntu 中运行。

\noindent{\color{red}操作失误，第三次实验作业不小心放到第四次里面了，现在已经移动到第三次实验了。}

\reportimage{images/hw06-02-agent-diagnosis.png}{Claude 分析除零错误并处理运行环境问题}

  \item {\large 由智能体修改代码并验证}

  Claude 在 \texttt{divide} 函数中增加了除数判断。\texttt{b == 0} 时返回提示字符串，
  其他情况仍执行原来的除法：

\begin{lstlisting}[language=Python]
def divide(a, b):
    if b == 0:
        return "Error: division by zero!"
    return a / b
\end{lstlisting}

  修改后重新运行，终端依次输出 5.0 和 \texttt{Error: division by zero!}，
  没有再出现异常回溯。

\noindent{\color{red}操作失误，第三次实验作业不小心放到第四次里面了，现在已经移动到第三次实验了。}

\reportimage{images/hw06-03-agent-fix.png}{Claude 修改 divide.py 并运行验证}

  \item {\large 检查最终代码与运行结果}

  退出智能体后，我用 \texttt{cat divide.py} 查看最终代码，再手动执行
  \texttt{python3 divide.py}。两次输出与智能体的运行结果一致，除零时程序不会再崩溃。

\noindent{\color{red}操作失误，第三次实验作业不小心放到第四次里面了，现在已经移动到第三次实验了。}

\reportimage{images/hw06-04-final-check.png}{检查修复后的代码并再次运行程序}
\end{enumerate}

\newpage
\section{第 7 题 编写一个最简单的 README}

\begin{enumerate}
  \item {\large 创建简单 Python 项目}

  在 \texttt{assignment-3/homework} 下创建 \texttt{07} 目录，并在其中建立
  \texttt{hello.py}。该程序只包含一条输出语句：

\begin{lstlisting}[language=Python]
print("Hello, World!")
\end{lstlisting}

  \item {\large 编写 README.md}

  在项目目录中创建 \texttt{README.md}，写清项目用途、文件组成、运行命令和输出结果：

\begin{lstlisting}[basicstyle=\ttfamily\footnotesize]
# Hello Project

这是一个简单的 Python 程序，用于输出 `Hello, World!`。

## 文件说明
- `hello.py`：Python 程序文件。
- `README.md`：项目说明文件。

## 运行方法
python3 hello.py

## 运行结果
Hello, World!
\end{lstlisting}

  按照 README 中的命令即可运行程序，不需要再打开源码查找入口。

\noindent{\color{red}操作失误，第三次实验作业不小心放到第四次里面了，现在已经移动到第三次实验了。}

\reportimage{images/hw07-readme.png}{创建 hello.py 并编写 README.md}
\end{enumerate}

\newpage
\section{第 8 题 编写一个规范的 Bug 报告}

\begin{enumerate}
  \item {\large 创建并运行错误程序}

  在 \texttt{assignment-3/homework} 下创建 \texttt{08} 目录，并编写
  \texttt{bug.py}。程序建立一个只有三个元素的列表，却尝试访问下标为 5 的元素：

\begin{lstlisting}[language=Python]
numbers = [1, 2, 3]
print(numbers[5])
\end{lstlisting}

  运行 \texttt{python3 bug.py} 后，程序在第 2 行触发
  \texttt{IndexError: list index out of range}。终端同时给出了文件路径和出错代码行，
  这些信息都要写进 Bug 报告。

  \item {\large 编写 BUG\_REPORT.md}

  创建 \texttt{BUG\_REPORT.md}，按照题目要求依次记录环境、期望结果、实际结果和复现步骤：

\begin{lstlisting}[basicstyle=\ttfamily\footnotesize]
# Bug Report

## 环境
- 系统：WSL2 Ubuntu
- Python：Python 3

## 期望结果
程序正常输出列表中的数字。

## 实际结果
程序出现 IndexError。

## 复现步骤
1. 进入程序目录。
2. 运行 `python3 bug.py`。
3. 程序出现 IndexError。
\end{lstlisting}

  复现步骤只保留进入目录和运行程序这两个必要操作。期望结果与实际结果分开写，
  查看报告时可以直接看出问题所在。

\noindent{\color{red}操作失误，第三次实验作业不小心放到第四次里面了，现在已经移动到第三次实验了。}

\reportimage{images/hw08-bug-report.png}{复现 IndexError 并创建规范的 Bug 报告}
\end{enumerate}

\newpage
\section{第 9 题 手动规范 Python 代码格式并提高可读性}

\begin{enumerate}
  \item {\large 创建并查看待整理代码}

  在 \texttt{assignment-3/homework} 下创建 \texttt{09} 目录，并编写
  \texttt{messy.py}。原代码能够完成加法运算，但运算符和函数参数之间缺少空格：

\begin{lstlisting}[language=Python]
def add(a, b):
    return a+b


print(add(2,3))
\end{lstlisting}

  代码虽然能运行，但空格使用不统一，看起来比较挤。

  \item {\large 手动规范代码格式}

  重新编辑 \texttt{messy.py}，在加法运算符两侧和函数调用的逗号后增加空格，并在函数定义与
  顶层调用之间保留两个空行。整理后的代码如下：

\begin{lstlisting}[language=Python]
def add(a, b):
    return a + b


print(add(2, 3))
\end{lstlisting}

  这次只改格式，没有动 \texttt{add} 的计算逻辑。执行 \texttt{python3 messy.py}
  后仍输出 5。

\noindent{\color{red}操作失误，第三次实验作业不小心放到第四次里面了，现在已经移动到第三次实验了。}

\reportimage{images/hw09-format-code.png}{手动规范 messy.py 的代码格式并运行}
\end{enumerate}

\newpage
\section{第 10 题 使用 pytest 为简单 Python 函数编写并运行单元测试}

\begin{enumerate}
  \item {\large 编写待测试函数}

  在 \texttt{assignment-3/homework} 下创建 \texttt{10} 目录，并在
  \texttt{calculator.py} 中编写加法函数：

\begin{lstlisting}[language=Python]
def add(a, b):
    return a + b
\end{lstlisting}

  \item {\large 编写单元测试}

  创建 \texttt{test\_calculator.py}，从 \texttt{calculator} 模块导入 \texttt{add}，
  再通过断言检查 \texttt{add(2, 3)} 的返回值是否等于 5：

\begin{lstlisting}[language=Python]
from calculator import add


def test_add():
    assert add(2, 3) == 5
\end{lstlisting}

  测试文件和测试函数都以 \texttt{test\_} 开头，Pytest 可以自动找到它。

  \item {\large 运行测试并查看结果}

  在项目目录执行安静模式测试：

\begin{lstlisting}[language=bash]
pytest -q
\end{lstlisting}

  终端显示 \texttt{1 passed in 0.01s}。运行 \texttt{ls} 后可以看到
  \texttt{calculator.py}、\texttt{test\_calculator.py}，以及测试时生成的
  \texttt{\_\_pycache\_\_} 目录。

\noindent{\color{red}操作失误，第三次实验作业不小心放到第四次里面了，现在已经移动到第三次实验了。}

\reportimage{images/hw10-pytest.png}{使用 Pytest 运行加法函数的单元测试}
\end{enumerate}

\color{red}\chapter{解题感悟}\color{black}
\section{课上练习}
\begin{enumerate}
  \item {\large 第 9 题是我第一次把一个小项目从源码构建成 Wheel，再放到新虚拟环境中安装。
  特意离开源码目录运行 \texttt{sdt-greet} 后仍能得到正确输出，这一步比只看“安装成功”更直观。}
  \item {\large 第 10 题里最有用的是先保留失败测试。它把“空白用户名处理不对”变成了明确的
  \texttt{1 failed}。智能体修好后，我又看了一遍差异并重跑测试，能够确定它没有顺手改动别的文件。}
  \item {\large 第 11 题让我发现，“程序无法处理”这种说法几乎不能帮助别人定位问题。
  把环境、命令、期望和实际结果补齐后，同一个问题就清楚多了；评审意见也应分清必须修改和建议改进。}
  \item {\large 第 12 题把 PyTorch 训练循环的顺序串了起来：先清空梯度，再反向传播，最后更新参数。
  训练结束不能只看程序是否跑完，还要检查 loss、weight 和 bias 是否接近题目给出的目标。}
\end{enumerate}

\newpage
\section{课后习题}
\begin{enumerate}
  \item {\large 第 1 题实际做下来，我才把“写配置、建虚拟环境、安装本地包、生成锁文件”这几步
  连在一起。以前看到 \texttt{pyproject.toml} 和锁文件时只知道它们与依赖有关，这次通过
  \texttt{pip list} 和 \texttt{uv.lock} 对照，能看懂包名、版本和来源记录在哪里了。}
  \item {\large 第 6 题给我的印象最深。除零错误本身不难，但 Claude 一开始受 Windows 和 WSL
  路径影响，连续两次没有把程序跑起来。这让我意识到，智能体给出修改并不等于任务已经完成，
  还要看它是否在正确的环境中真正运行过。最后我又手动执行了一次，心里才有底。}
  \item {\large 第 10 题让我体会到单元测试的作用。以前修改小函数时通常只看一次输出，
  现在可以把 \texttt{add(2, 3) == 5} 写成断言，以后代码发生变化时直接运行 Pytest，
  很快就能知道原来的功能有没有被改坏。}
\end{enumerate}

\color{red}\chapter{上传GitHub}\color{black}

\section{编写并更新 .gitignore}

提交第三次实验前，我在课程仓库根目录继续更新 \texttt{.gitignore}。本次实验会产生虚拟环境、Python 构建目录、\texttt{egg-info} 元数据和测试缓存，如果全部上传，不但文件数量多，也会把本机环境带进仓库。因此，在原有规则后补充了第三次实验相关内容：

\begin{lstlisting}
# 第三次实验的虚拟环境和 Python 构建中间产物
/assignment-3/q09-test/
/assignment-3/**/build/
/assignment-3/**/*.egg-info/

# q10 从 q09 复制的旧构建包；q09/dist 保留为第 9 题实验成果
/assignment-3/q10/dist/

# 第三次实验报告的排版检查图片和本地备份
/assignment-3/实验报告/tmp/
/assignment-3/实验报告/*.bak
.venv/
\end{lstlisting}

其中，\texttt{q09-test} 和 \texttt{.venv} 都是可以重新创建的虚拟环境；\texttt{build}、\texttt{egg-info} 和 \texttt{q10/dist} 属于构建过程中产生的文件。实验报告的临时检查图片和备份文件也没有提交。这样处理后，仓库中保留源码、测试、作业截图和正式报告，提交内容更清楚。

\section{GitHub 仓库链接}

本课程作业保存在下面的公开仓库中，默认分支为 \texttt{main}：

\begin{center}
\url{https://github.com/xiaoyang1088/system_development_tools}
\end{center}

仓库用户名为 \texttt{xiaoyang1088}，仓库名为 \texttt{system\_development\_tools}。第三次实验位于 \texttt{assignment-3/} 目录，其中包括课堂练习 \texttt{q09}--\texttt{q12}、课后练习 \texttt{homework/}、实验截图和实验报告。

\section{commit 提交记录}

第三次实验分三次提交。2026 年 9 月 14 日先提交了针对第三次实验补充的 \texttt{.gitignore}，对应提交号 \texttt{246d4c9}；9 月 15 日提交课堂练习 \texttt{q09}--\texttt{q12}，对应提交号 \texttt{2d3b9b7}；随后提交课后练习 \texttt{homework}，对应提交号 \texttt{b861116}。分开提交后，每次修改的范围比较明确，后续查看或回退时也更方便。

\reportimage{images/github-commit-history.png}{第三次实验作业的 GitHub commit 提交记录}

\end{document}
